\documentclass[12pt]{article}

\usepackage[a4paper,margin=2.5cm]{geometry}

\usepackage[onehalfspacing]{setspace}

\usepackage{amsmath}
\usepackage{amssymb}
\usepackage{xcolor}
\usepackage{mathtools}
\usepackage{amsthm}
\newtheorem{definition}{Definition}
\newtheorem{example}{Example}

\setlength{\parindent}{0pt}

\begin{document}

\begin{titlepage}
    \centering

    \vspace*{\fill}

    {\Huge Basis and Dimension I}

    \vspace{1cm}

    {\Large Jin-Ren Ryan Wang\par}

    \vspace{0.5cm}

    {\large \today\par}

    \vspace*{\fill}
\end{titlepage}

\tableofcontents

\newpage

\section{Linear combination}

\subsection{Definition}
\textbf{Definition 2.2.1.}

Let $S$ = \{$\mathbf{v_1}, \mathbf{v_2},\ldots, \mathbf{v_m}$\} be a subset of $V$. Then the expression
\begin{center}
    $a_1\mathbf{v_1} + a_2\mathbf{v_2} + \ldots + a_m\mathbf{v_m}$, with $a_i \in F$
\end{center}


is called a $linear$ $combination$ of $S$.

\subsection{Example}

\textbf{1.}
(2, 1, 0) is A linear combination of $S$ = \{(1, 2, 3), (4, 5, 6)\} because
\[
    (2, 1, 0) 
    = 
    -2 \cdot (1, 2, 3) + 1 \cdot (4, 5, 6).
\]

\textbf{2.}
$f(x)$ = $x^{3} + 2x + 1$ in $\mathbb{R}[x]$ is a linear combination of $S$ = $\{1,x,x^{2},x^{3}\}$ because
\[
    f(x) = 1 \cdot 1 + 2 \cdot x + 0 \cdot x^{2} + 1 \cdot x^{3}.
\]

\textbf{3. (polynomial)} 
Show that $\{2x+1,\;x^2+3\}$ is linearly independent in $\mathbb{R}[x]_{\le 2}$.\\
Suppose
\[
    a_1(2x+1)+a_2(x^2+3)=0. 
\]
\[
    \implies a_2x^2+2a_1x+(a_1+3a_2)=0. 
\]
By comparing coefficients yields
\[
a_2=0,\quad 2a_1=0,\quad a_1+3a_2=0.\implies a_1=a_2=0
\]
Therefore
$\{2x+1,\;x^2+3\}$ is linearly independent.    

\newpage

\section{Linear dependence and independence}

\subsection{Definition}
\textbf{Definition 2.2.2.}
\begin{enumerate}
    \item A \textcolor{red}{finite} subset $S$ = $\{\mathbf{v_1},\mathbf{v_2},\ldots,\mathbf{v_m}\}$ of $V$ is called $linearly$ $independent$ if
    \[
        a_1\mathbf{v_1} + a_2\mathbf{v_2} + \ldots + a_m\mathbf{v_m} = \mathbf{0} \implies a_1 = a_2 = \ldots = a_m = 0.
    \]   In other words, the only way to obtain $\mathbf{0}$ as a linear combination of $S$ is by taking all scalars to 0.

    \item More generally, an \textcolor{red}{infinite} subset $S'$ of $V$ is said to be $linearly$ $independent$ if every finite subset $S \subseteq S'$ is $linearly$ $independent$.

    \item A subset $S''$ is called $linearly$ $dependent$ if it is $not$ $linearly$ $independent$.

\end{enumerate}

\subsection{Example}

Prove that, in $\mathbb{R}^3$,

\begin{enumerate}
    \item $S_1$ = \{(1, 2, 3), (4, 5, 6)\} is linearly independent:
    \begin{proof}
        Suppose
        \[
        a_1(1,2,3) + a_2(4,5,6) = (0,0,0).
        \]
        Then
        \[
        \begin{cases}
        a_1 + 4a_2 = 0,\\
        2a_1 + 5a_2 = 0,\\
        3a_1 + 6a_2 = 0.
        \end{cases}
        \]
        From the first two equations,
        \[
        2(a_1+4a_2)-(2a_1+5a_2) = 3a_2 = 0.
        \]
        Hence $a_2 = 0$, and therefore $a_1 = 0$, so the trivial solution is the only solution.\\
        $\implies$ $\{(1,2,3),(4,5,6)\}$ is linearly independent.
    \end{proof}
    \item $S_2$ = \{(1, 2, 3), (4, 5, 6), (7, 8, 9)\} is linearly dependent.  
        \begin{proof}
            As noted previously, 
            \[
                1 \cdot (1, 2, 3) + (-2) \cdot (4, 5, 6) + 1 \cdot (7, 8, 9) = (0, 0, 0)
            \] By def, $S_2$ is not linearly independent.
        \end{proof}
\end{enumerate}

\subsection{Theorem}

\textbf{Theorem 2.2.1.} 

$\{\mathbf{v_1}, \mathbf{v_2}, \ldots, \mathbf{v_m}\}$ is linearly dependent iff. either $\mathbf{v_1}$ = $\mathbf{0}$ or
for some r, $\mathbf{v_r}$ is a linear combination of  $\mathbf{v_1}, \mathbf{v_2}, \ldots, \mathbf{v_{r - 1}}$
\begin{proof}

\textcolor{red}{($\Rightarrow$)} Given $S=\{\mathbf{v_1},\dots,\mathbf{v_m}\}$ is NOT linearly dependent.

\textbf{Case 1.} $\mathbf{v_1}=\mathbf{0}$.$\quad$Then there is nothing to prove.

\textbf{Case 2.} $\mathbf{v_1}\neq \mathbf{0}$.

Since $S$ is linearly dependent, there exist scalars
$a_1, \ldots, a_m$, not all zero, s.t.

\[
a_1\mathbf{v_1} + \cdots + a_m\mathbf{v_m}=\mathbf{0}.
\]

Let $r$ be the largest index such that

\begin{center}
    $a_r\neq 0.$ $(a_{r+1} = a_{r+2} = \ldots = a_m = 0)$
\end{center}

Hence $(*)$ becomes
\[
a_1\mathbf{v_1} + a_2\mathbf{v_2} + \cdots + a_r\mathbf{v_r} = 0.
\]

Since $a_r\neq 0$, we may solve for $\mathbf{v_r}$:
\[
\mathbf{v_r}
=
-\frac{a_1}{a_r}\mathbf{v_1}
-\frac{a_2}{a_r}\mathbf{v_2}
-\cdots
-\frac{a_{r-1}}{a_r}\mathbf{v_{r-1}}.
\]

Therefore $\mathbf{v_r}$ is a linear combination of
\[
\{\mathbf{v_1},\dots,\mathbf{v_{r-1}}\}.
\]

\textcolor{red}{($\Leftarrow$)} 

\textbf{Case 1.} Suppose $\mathbf{v_1} = \mathbf{0}$\\
$\Rightarrow$ $1\mathbf{v_1} + 0\mathbf{v_2} + 0\mathbf{v_3} + \ldots = 0\mathbf{v_m}$ = $\mathbf{0}$ \\
$\Rightarrow$ \textcolor{red}{The coefficients are NOT all zero.}

\textbf{Case 2.} $\mathbf{v_1}\neq 0$.

Assume that there  $\exists r$ s.t.
\[
\mathbf{v_r} = a_1\mathbf{v_1} + a_2\mathbf{v_2} + \cdots + a_{r-1}\mathbf{v_{r-1}}.
\]

Then
\[
a_1\mathbf{v_1} + a_2\mathbf{v_2} + \cdots + a_{r-1}\mathbf{v_{r-1}} - \mathbf{v_r} = 0.
\]

Since the coefficient of $\mathbf{v_r}$ is $-1\neq 0$,
the coefficients are not all zero.

Therefore there exists a nontrivial linear relation among the vectors in $S$.

Hence $S$ is linearly dependent.

\end{proof}

\section{Spanning sets}

\subsection{Definition}

\textbf{Definition 2.3.1. (Linear span)}

\begin{enumerate}
    \item Let $S = \{\mathbf{v_1}, \mathbf{v_2}, \ldots, \mathbf{v_m}\}$ be a finite subset of $V$. Then define the (linear) span of $S$ by
    \[
        \operatorname{span}(S) \coloneq \{a_1\mathbf{v_1} + a_2\mathbf{v_2} + \ldots + a_m\mathbf{v_m}\} : a_i \in F
    \]
    the set of all linear combination of $S$.\\
    i.e. span($S$) = \{all linear combination of $S$\}
    \item More generally, if $S$ is an infinite subset of $V$ (i.e. $S$ contains infinitely many elements). then we define
    \[
        \operatorname{span}(S)
        =
        \bigcup_{S'\subset S} \operatorname{span}(S')    
    \]
    where $S'$ runs through all finite subsets of $S$.
\end{enumerate}

\textbf{Definition 2.3.2. (Span as a verb)}

If $V$ = $\operatorname{span}(S)$ (i.e. every $\mathbf{v} \in V$ can be written as a linear combination
of members in $S$.), we say that $S$ spans $V$ or $S$ is a spanning set of $V$.

\subsection{Remark}

By Convention, we have 
\fcolorbox{red}{white}{$\operatorname{span}(\emptyset) = \{\mathbf{0}\}$}.
\qquad
$\leftarrow$
\textcolor{red}{empty sum is defined to be $0$}

\subsection{Theorem}
\textbf{Therorem 2.3.1.}

$\operatorname{span}(S)$ is a subspace of $V$.

\begin{proof}
This is straight-forward, so we do it "verbally":

\begin{enumerate}
    \item \textbf{(Closed under addition)}
    
    The sum of any two linear combinations of $S$ is still a linear combination of $S$ (by adding up the coefficients).

    \item \textbf{(Closed under scalar multiplication)}
    
    A scalar multiple of a linear combination of $S$ is still a linear combination of $S$ (by scaling the coefficients).

    \item
    
    If $S$ is non-empty, then $\operatorname{span}(S)$ is also non-empty. If $S$ is empty, then
    \[
        \operatorname{span}(S)=\{0\},
    \]
    which is also non-empty.
\end{enumerate}
\end{proof}

\subsection{Example}

\textbf{1.}

\bigskip

Let $S$ = \{(1, 0, 0), (0, 1, 0), (0, 0, 1)\} $\subseteq$ $\mathbb{R}^{3}$
\begin{align*}
\operatorname{span}(S) 
&\coloneq
\{a(1, 0, 0) + b(0, 1, 0) + c(0, 0, 1) : a, b, c \in \mathbb{R}\}\\
&= 
\{(a, b, c) : a, b, c \in \mathbb{R}^{3}\}\\    
&= 
\mathbb{R}^{3}
\end{align*}

\textbf{2.}

\bigskip

$S'$ = \{(1, 0, 0), (0, 1, 0), (0, 0, 1), (1, 1, 1)\} also spans $\mathbb{R}^{3}$ because
More concretely, for any (a, b, c) $\in \mathbb{R}$,
\[
\underbrace{(a, b, c) = a(1, 0, 0) + b(0, 1, 0) + c(0, 0, 1) + 0(1, 1, 1)
}_{\text{linear combination of $S'$}}
\]

\newpage

\section{Basis}

\subsection{Definition}

\textbf{Definition 2.4.1.}

A subset $S$ of $V$ is called a \emph{basis} of $V$ if

\begin{description}
    \item[(1)] $S$ is linearly independent.
    \item[(2)] $S$ spans $V$, i.e.\ every $\mathbf{v}\in V$ can be written as a linear combination of $S$.
\end{description}

Moreover, if $S$ can be taken to be a finite set, then we say that $V$ is a
\emph{finite-dimensional} vector space. Otherwise, we say that it is
\emph{infinite-dimensional}.

\bigskip

\textbf{Definition 2.4.2.}

Fix an order of elements in a basis $S$. In this case, we call
$(a_1,a_2,\ldots,a_m)$ the \emph{coordinates} of $\mathbf{v}$ with respect to the basis
$S$ and write

\[
[\mathbf{v}]_S=(a_1,a_2,\ldots,a_m).
\]

\subsection{Therorem}

\textbf{Theorem 2.4.1.}

Suppose
\[
S=\{\mathbf{v_1,v_2,\ldots,v_m}\}
\]
is a basis of $V$. Then for every $v\in V$, there exists a unique collection of scalars
$a_1,a_2,\ldots,a_m\in F$ such that

\[
\mathbf{v}=a_1\mathbf{v_1}+a_2\mathbf{v_2}+\cdots+a_m\mathbf{v_m}.
\]

\begin{proof}
Since $S$ is a basis of $V$, it spans $V$. Therefore every $v\in V$ can be
written as a linear combination of the vectors in $S$.

To prove uniqueness, suppose
\[
\mathbf{v}=a_1\mathbf{v_1}+a_2\mathbf{v_2}+\cdots+a_m\mathbf{v_m}
\]

and
\[
\mathbf{v}=b_1\mathbf{v_1}+b_2\mathbf{v_2}+\cdots+b_m\mathbf{v_m},
\]

where $a_i,b_i\in F$.

Subtracting the two equations, we obtain
\[
\mathbf{0}=(a_1-b_1)\mathbf{v_1}+(a_2-b_2)\mathbf{v_2}+\cdots+(a_m-b_m)\mathbf{v_m}.
\]

Since $S$ is a basis, it is linearly independent. Hence
\[
(a_1-b_1)=(a_2-b_2)=\cdots=(a_m-b_m)=0.
\]

Therefore
\[
a_1=b_1,\quad a_2=b_2,\quad \ldots,\quad a_m=b_m.
\]

Thus the representation of $v$ is unique.
\end{proof}

\subsection{Example}

\textbf{1. Non-uniqueness of a basis}

\bigskip

The set
\[
S=\{(1,0,0),(1,1,0),(1,1,1)\}
\]

is also a basis for $\mathbb{R}^3$ because it is linearly independent and for every
$(a,b,c)\in\mathbb{R}^3$,

\[
(a,b,c)
=
(a-b)(1,0,0)
+
(b-c)(1,1,0)
+
c(1,1,1).
\]

Hence every vector in $\mathbb{R}^3$ is a linear combination of vectors in $S$.

Moreover, since $S$ is finite, $\mathbb{R}^3$ is finite-dimensional.

\bigskip

\textbf{2. Infinite-dimensional example}

\bigskip

Consider $\mathbb{R}[x]$.

We claim that $\mathbb{R}[x]$ cannot be spanned by any finite set of polynomials.

\begin{proof}
Suppose
\[
S=\{f_1,f_2,\ldots,f_m\}
\]

spans $\mathbb{R}[x]$.

Let
\[
N=\max\{\deg(f_1),\deg(f_2),\ldots,\deg(f_m)\}.
\]

Then every polynomial in $\operatorname{span}(S)$ is a linear combination of
$f_1,f_2,\ldots,f_m$, and therefore has degree at most $N$.

However,
\[
x^{N+1}
\]

has degree $N+1$, so
\[
x^{N+1}\notin \operatorname{span}(S).
\]

This contradicts the assumption that $S$ spans $\mathbb{R}[x]$.

\bigskip

Therefore, $\mathbb{R}[x]$ is infinite-dimensional.
\end{proof}

\bigskip

\textbf{3.}

\bigskip

As observed previously, for
\[
\mathbf{v}=(a,b,c)\in\mathbb{R}^3,
\]

we have

\begin{itemize}

\item with respect to the standard basis
\[
S_1=\{(1,0,0),(0,1,0),(0,0,1)\},
\]

we have
\[
[\mathbf{v}]_{S_1}=(a,b,c).
\]

\item (\emph{Order matters}) with respect to the reordered standard basis
\[
S_1'=\{(0,1,0),(0,0,1),(1,0,0)\},
\]

we have
\[
[\mathbf{v}]_{S_1'}=(b,c,a).
\]

\item with respect to the basis
\[
S_2=\{(1,0,0),(1,1,0),(1,1,1)\},
\]

we have
\[
[\mathbf{v}]_{S_2}=(a-b,b-c,c).
\]

\end{itemize}

\subsubsection{Remark}
In particular, the same vector $\mathbf{v}$ may have different coordinates with respect to a different choices of basis.

\newpage

\section{Dimension}

\subsection{Definition}

\textbf{Definition 2.5.1}

Suppose $V$ is a finite-dimensional vector space over $F$.

Define the dimension of $V$ to be the size of a (and hence any) basis of $V$:

\[
\dim_F(V)=|S|
\]

where $S$ is a basis of $V$.


\subsection{Theorem}

\textbf{Theorem 2.5.1.(Exchange Theorem / Replacement Theorem)}

Suppose

\begin{itemize}
    \item $S_1=\{\mathbf{v_1,v_2,\ldots,v_n}\}$ is a spanning set of $V$;

    \item $S_2=\{\mathbf{w_1,w_2,\ldots,w_m}\}$ is linearly independent.
\end{itemize}

Then
\[
m\le n.
\]

\begin{proof}
    This proof will be provided after develping \emph{sifiting method} theorem in\\
    other chapter.
\end{proof}

\textbf{Theorem 2.5.2. (Basis Theorem)}

Suppose $V$ is a finite-dimensional vector space over $F$.
If $S_1$ and $S_2$ are both bases of $V$, then $S_1$ and $S_2$
contain the same number of elements.

\begin{proof}
Let $S_1$ and $S_2$ be two bases of $V$.

\begin{itemize}
    \item Since $S_1$ is a basis, it is linearly independent.
    
    Since $S_2$ is a basis, it is a spanning set.
    
    By the Exchange Theorem,
    \[
    |S_1|\le |S_2|.
    \]

    \item Since $S_2$ is a basis, it is linearly independent.

    Since $S_1$ is a basis, it is a spanning set.
    
    By the Exchange Theorem,
    \[
    |S_2|\le |S_1|.
    \]

\end{itemize}

Hence,

\[
|S_1|=|S_2|.
\]

Therefore every basis of $V$ has the same cardinality.
\end{proof}

\subsection{Example}

\textbf{1.}

\bigskip

Since
\[
S=\{(1,0,0),(0,1,0),(0,0,1)\}
\]

is a basis of $\mathbb{R}^3$ (as a vector space over $\mathbb{R}$), we conclude that
\[
\dim_{\mathbb{R}}(\mathbb{R}^3)
=
|S|
=
3.
\]

\end{document}